\documentclass[11pt,a4paper]{article}
\usepackage[T1]{fontenc}
\usepackage[utf8]{inputenc}
\usepackage{amsmath,amssymb,xcolor}
\usepackage[margin=25mm]{geometry}
\usepackage[colorlinks=true,urlcolor=blue]{hyperref}
\definecolor{nexorablue}{HTML}{2459B5}
\setlength{\parindent}{0pt}
\setlength{\parskip}{7pt}
\setlength{\emergencystretch}{2em}
\renewcommand{\familydefault}{\sfdefault}
\begin{document}

{\small\bfseries\color{nexorablue} NEXORA CPGE}\par
{\small MP / NEX-STR-02 / Énoncé}\par
{\LARGE\bfseries Ordre et image d’un sous-groupe cyclique}\par
\vspace{0.25em}\hrule\vspace{0.6em}

Dans le groupe additif $\mathbb Z/30\mathbb Z$, déterminer l’ordre de $\overline{12}$, puis décrire le sous-groupe qu’il engendre. Préciser le noyau et l’image de $n\mapsto\overline{12n}$ sur $\mathbb Z$.


\end{document}
